\scanpage{45}
\mspage{1}
% archive foliation: 44
1. Opération de groupes $\simeq\mathbb D_2$ sur la droite projective, et configuration octaédrale.

Soit $V$ un groupe qui est un espace vectoriel de rang $2$ sur $\mathbb F_2$, et considérons \unsure{un|une} \deleted{droite pro} fibré en droites projectives $X$ sur un schéma de base $S$. On s'intéresse aux opérations, fidèles sur chaque fibre, de $V$ sur $X$, i.e. aux monomorphismes
\[
(1)\qquad V_S\lhook\joinrel\longrightarrow\underline{\operatorname{Aut}}_S(X)=G
\]
Choisissons une base $e_1,e_2$ de $V$ sur $\mathbb F_2$, alors la donnée d'un hom
\[
V_S\longrightarrow\underline{\operatorname{Aut}}_S(X)\qquad\text{i.e.}\quad V\longrightarrow\operatorname{Aut}_S(X)
\]
qui soit mono sur $\mathbb F_2e_1$ et sur $\mathbb F_2e_2$, équivaut à celle de deux \inserted{\unsure{autom.} d'ordre deux de $X$, ou encore de} ``axes'' ($=$ diviseurs étales de degré $2$) $\Delta_1,\Delta_2$ sur $X$, la condition de commutation étant que $\Delta_2$ est stable par $\sigma_1$, ou encore $\Delta_1$ \unsure{stble|stable} par $\sigma_2$.
\margin{on est en car. $\ne2$\quad $\Delta_1=X^{\sigma_1}$\quad $\Delta_2=X^{\sigma_2}$}
\deleted{Pour $\sigma_1$ donné, i.e. $\Delta_1$ donné, la donnée de $\sigma_2$ équivaut donc à la donnée d'une section de $X/\{1,\sigma_1\}\longrightarrow$}
\inserted{La condition que (1) soit un mono signifie \unsure{aussi} que l'on \unsure{a :}}
\[
(2)\qquad\Delta_1\cap\Delta_2=\varnothing\qquad\text{i.e.}\quad\Delta_2\subset X\unsure{-|\setminus}\Delta_1
\]
et pour $\sigma_1$ \inserted{i.e. $\Delta_1$} donné, la donnée de \deleted{$\unsure{\Delta_2}$} $\sigma_2$ i.e. de $\Delta_2$ équivaut à celle d'une section de \deleted{$\unsure{X\unsure{-|\setminus}\Delta_1}$ \unsure{sur}} $(X\unsure{-|\setminus}\Delta_1)/\{1,\sigma_1\}$ sur $S$ ($\Delta_2$ étant alors l'image inverse de cette section dans $X\unsure{-|\setminus}\Delta_1$). Considérons le troisième larron
\[
(3)\qquad e_3=e_1+e_2
\]
il lui correspond un axe $\Delta_3$, et $\Delta_1,\Delta_2,\Delta_3$ sont mutuellement disjoints.
